% Excerpt of Appendix G of "From Correlations to Geometry" (working copy of 17 November 2026).
% A fragment, not a standalone document; it uses the macros and labels of the paper.
% The typeset pages are in appendix_G22_G37.pdf.

% Lemmas G.26, G.27, Theorem G.28, Corollaries G.29-G.31, Lemma G.32, Theorem G.33, Corollaries G.34, G.35,
% Remarks G.36, G.37.

\paragraph{The kernel of the sampled walk} The remaining statements of this subsection concern the walk itself and not only its trace. For $t>0$ let $\gamma_t(v)=(2\pi t)^{-d/2}e^{-|v|^2/2t}$ on $\R^d$ and $p^{(t)}(w)=\sum_{n\in\mathbb Z^d}\gamma_t(w+n)$ on $\M=\R^d/\mathbb Z^d$, the heat kernel of $\M$ at the time $t/2$. Thus $p=p^{(\xi^2)}$, $p_m=p^{*m}=p^{(m\xi^2)}$, $M=p(0)$, $m_m=p_m(0)$, and $(\mathcal Pf)(x)=\int_\M p(x-y)f(y)\,\dd y$. The points $x_1,\dots,x_N$ are independent and uniform on $\M$, $\tilde P_{bj}=p(x_b-x_j)/N$ (self-loops included), $\epsilon_b=\sum_j\tilde P_{bj}-1$, $\delta=\max_b|\epsilon_b|$, $u_j=(1+\epsilon_j)^{-1}$, $r_b=\sum_j\tilde P_{bj}u_j$ and $P_{bj}=\tilde P_{bj}u_j/r_b$. For $1\le i\le N$ and $m\ge1$,
\[
h^{(i)}_m(b)=\sum_jP_{bj}\,p_m(x_j-x_i)-p_{m+1}(x_b-x_i).
\]

\begin{lemma}[Facts on the periodised Gaussian]\label{lem:pgauss}
There is a constant $C_d\ge1$ such that for all $t>0$, $0<a\le\min(t,1)$ and $u,z\in\M$:
\begin{enumerate}[label=\textup{(P\arabic*)}]
\item $p^{(a)}\le(t/a)^{d/2}\,p^{(t)}$ pointwise;
\item $\sup p^{(t)}=p^{(t)}(0)$ is non-increasing in $t$; $p^{(t)}(0)\le C_dt^{-d/2}$ for $t\le2$, $p^{(t)}(0)\le C_d$ for $t\ge2$, and $p^{(t)}\ge\frac12$ for $t\ge t_d$, where $t_d\ge1$ depends only on $d$;
\item $|\nabla p^{(t)}|\le2^{d/2}\,t^{-1/2}\,p^{(2t)}$ pointwise;
\item $p^{(a)}(u)\,p^{(t)}(z-u)\le C_d\,(2\pi a)^{-d/2}\,p^{(3t)}(z)$;
\item $\big(p^{(2t)}(0)\,p^{(3t)}(z)\big)^{1/2}\le C_d\,p^{(6t)}(z)$.
\end{enumerate}
\end{lemma}

\begin{proof}
(P1): $\gamma_a\le(t/a)^{d/2}\gamma_t$ for $a\le t$. (P2): $p^{(t)}(w)=\sum_{k\in\mathbb Z^d}e^{-2\pi^2|k|^2t}e^{2\pi ik\cdot w}$ has positive coefficients that decrease in $t$; $p^{(t)}(0)=(2\pi t)^{-d/2}\big(\sum_{j\in\mathbb Z}e^{-j^2/2t}\big)^d\le(2\pi t)^{-d/2}\big(1+2\sum_{j\ge1}e^{-j^2/4}\big)^d$ for $t\le2$; and $|p^{(t)}-1|\le\sum_{k\ne0}e^{-2\pi^2|k|^2t}=\big(1+2\sum_{j\ge1}e^{-2\pi^2j^2t}\big)^d-1\le\frac12$ for $t\ge t_d$. (P3): $\nabla\gamma_t(v)=-(v/t)\gamma_t(v)$ and $re^{-r^2/2}\le e^{-r^2/4}$ for $r\ge0$ give $|\nabla\gamma_t|\le t^{-1/2}2^{d/2}\gamma_{2t}$; sum over the lattice. (P4): for $v,v'\in\R^d$, $\frac{|v|^2}{2a}+\frac{|v'|^2}{2t}\ge\frac{|v|^2}{4a}+\frac{|v+v'|^2}{2(t+2a)}$, because $\frac{|v|^2}{4a}+\frac{|v'|^2}{2t}\ge\frac{|v+v'|^2}{4a+2t}$ by the Cauchy--Schwarz inequality. With $v=u+n$, $v'=z-u+n'$ and $n''=n+n'$,
\[
\begin{aligned}
p^{(a)}(u)\,p^{(t)}(z-u)&\le(2\pi a)^{-d/2}(2\pi t)^{-d/2}\sum_ne^{-|u+n|^2/4a}\ \sum_{n''}e^{-|z+n''|^2/2(t+2a)}\\
&\le C\,(2\pi a)^{-d/2}\Big(\frac{t+2a}{t}\Big)^{d/2}p^{(t+2a)}(z),
\end{aligned}
\]
since $\sum_ne^{-|u+n|^2/4a}\le(2+2\sum_{j\ge1}e^{-j^2/4})^d$ for $a\le1$; now use $t+2a\le3t$ and (P1). (P5): for $t\le1$, $\big(p^{(3t)}(z)\big)^{1/2}\le\sum_n\gamma_{3t}(z+n)^{1/2}=(24\pi t)^{d/4}p^{(6t)}(z)$ and $p^{(2t)}(0)\le C_d(2t)^{-d/2}$; for $1\le t\le t_d$ the same computation applies with $p^{(2t)}(0)\le C_d$ and $(24\pi t)^{d/4}\le C_d$; for $t\ge t_d$ the left-hand side is bounded and $p^{(6t)}\ge\frac12$, by (P2).
\end{proof}

\begin{lemma}[Consistency of the normalised walk]\label{lem:consistency}
Let $d\ge1$, $0<\xi\le1$, $N\ge3$, $M\log N\le N$, $\delta'=\sqrt{M\log N/N}$ and $6\delta'\le\frac18$. There is a constant $C$ depending only on $d$ such that, with probability at least $1-7N^{-3}$, the following hold for all $b,i\in\{1,\dots,N\}$ and all integers $1\le m\le\delta'^{-2}$:
\begin{enumerate}[label=\textup{(\roman*)}]
\item $\delta\le6\delta'$;
\item $\big|h^{(i)}_m(b)\big|\le C\,\delta'\,m^{-1/2}\,p_{6m}(x_b-x_i)$;
\item $\frac1N\sum_bp_{6m}(x_b-x_i)\le1+9\delta'$.
\end{enumerate}
In particular $\max_b|h^{(i)}_m(b)|\le C\delta'm^{-1/2}m_{6m}$ and $\frac1N\sum_b|h^{(i)}_m(b)|\le2C\delta'm^{-1/2}$.
\end{lemma}

\begin{proof}
We use Bernstein's inequality in the following form: if $X_1,\dots,X_n$ are independent and centred, $|X_j|\le B$ and $\mathbb EX_j^2\le v$, then for every $L>0$, with probability at least $1-2e^{-L}$, $|\sum_jX_j|\le\sqrt{2nvL}+\frac23BL$. Throughout $L=6\log N$, and $C$ denotes constants that depend only on $d$.

\emph{\textup{(i)}} is Step~1 of the proof of Theorem~\ref{thm:jointpoly} and fails with probability at most $2N^{-3}$.

\emph{\textup{(iii)}.} Fix $i$ and $m$. Given $x_i$, the variables $p_{6m}(x_b-x_i)-1$, $b\ne i$, are independent and centred, bounded by $\max(m_{6m},1)\le M$ and of variance at most $\int p_{6m}^2=m_{12m}\le M$, by (P2). With probability at least $1-2N^{-6}$ their sum is at most $\sqrt{12MN\log N}+4M\log N$ in absolute value, that is, at most $N(\sqrt{12}\,\delta'+4\delta'^2)$; the term $b=i$ contributes $m_{6m}/N\le\delta'^2$. Since $\delta'\le1$ this gives (iii), and the union over $i$ and $m\le\delta'^{-2}\le N$ costs at most $2N^{-4}$.

\emph{\textup{(ii)}.} Fix $b$, $i$, $m$, put $s=m\xi^2$, $z=x_b-x_i$, $g=p_m(\cdot-x_i)$ and $c_b=(\mathcal Pg)(x_b)=p_{m+1}(z)$. Since $\sum_jP_{bj}=1$,
\begin{equation}\label{eq:cons-split}
h^{(i)}_m(b)=\frac1{r_b}\,(S_1+S_2),\qquad S_1=\sum_j\tilde P_{bj}\big(g(x_j)-c_b\big),\qquad S_2=\sum_j\tilde P_{bj}(u_j-1)\big(g(x_j)-c_b\big).
\end{equation}
On the event (i), $\delta\le\frac18$, hence $1/r_b\le\frac{10}7$ and $|u_j-1|\le\frac87\delta$, so that
\[
\big|h^{(i)}_m(b)\big|\le\tfrac{10}7\big(|S_1|+\tfrac87\,\delta\,A\big),\qquad A=\frac1N\sum_jp(x_b-x_j)\,\big|g(x_j)-c_b\big| .
\]
For $j\notin\{b,i\}$ let $Y_j=p(x_b-x_j)\,(g(x_j)-c_b)$. Given $x_b$ and $x_i$ the $Y_j$ are independent and identically distributed, and $\mathbb EY_j=\int p(x_b-y)p_m(y-x_i)\,\dd y-c_b=0$.

\emph{A bound for $Y_j$.} By (P4) with $a=\xi^2$, $t=s$ and $(2\pi\xi^2)^{-d/2}\le M$, $p(x_b-y)\,p_m(y-x_i)\le CM\,p^{(3s)}(z)$ for all $y$, and $c_b=p^{(s+\xi^2)}(z)\le3^{d/2}p^{(3s)}(z)$ by (P1). Hence
\begin{equation}\label{eq:cons-B}
|Y_j|\le B:=C\,M\,p^{(3s)}(z).
\end{equation}
\emph{The variance.} For a function $f$ on $\M$, $\int_\M p(x_b-y)f(y)\,\dd y=\mathbb E\,f(x_b+\xi Z)$ with a standard Gaussian vector $Z$ in $\R^d$ ($f$ extended periodically). Hence $c_b=\mathbb E\,g(x_b+\xi Z)$ and, by the Gaussian Poincar\'e inequality, (P3), (P2) and $\mathcal Pp^{(2s)}=p^{(2s+\xi^2)}\le(\frac32)^{d/2}p^{(3s)}$,
\begin{align}
\mathbb EY_j^2&\le M\int p(x_b-y)\,(g(y)-c_b)^2\,\dd y=M\operatorname{Var}g(x_b+\xi Z)\le M\,\xi^2\,\big(\mathcal P|\nabla p_m|^2\big)(z)\notag\\
&\le M\,\xi^2\,\frac{2^d}{s}\,p^{(2s)}(0)\,\big(\mathcal Pp^{(2s)}\big)(z)\le v:=\frac{C\,M}{m}\,p^{(2s)}(0)\,p^{(3s)}(z),\label{eq:cons-v}
\end{align}
and $\mathbb E|Y_j|\le(v/M)^{1/2}$ by the Cauchy--Schwarz inequality.

\emph{Bernstein.} With probability at least $1-2N^{-6}$, by \eqref{eq:cons-B}, \eqref{eq:cons-v}, (P1) and (P5),
\[
\begin{aligned}
\Big|\frac1N\sum_{j\notin\{b,i\}}Y_j\Big|&\le\sqrt{\frac{12\,v\log N}{N}}+\frac{4B\log N}{N}\le C\,\delta'\,m^{-1/2}\big(p^{(2s)}(0)\,p^{(3s)}(z)\big)^{1/2}+C\,\delta'^2\,p^{(3s)}(z)\\
&\le C\,\delta'\,m^{-1/2}\,p^{(6s)}(z),
\end{aligned}
\]
where the last step uses $\delta'\le m^{-1/2}$. The same bound holds, with the same probability, for $\frac1N\sum_{j\notin\{b,i\}}(|Y_j|-\mathbb E|Y_j|)$, because $\big||Y_j|-\mathbb E|Y_j|\big|\le B$ and $\operatorname{Var}|Y_j|\le v$.

\emph{The two exceptional terms.} The term $j=b$ of $S_1$ and of $A$ is at most $\frac MN\,(p_m(z)+p_{m+1}(z))\le C\,\frac MN\,p^{(3s)}(z)$ in absolute value, by (P1). For $b\ne i$ the term $j=i$ is at most $\frac1N\,(p(z)\,p_m(0)+M\,c_b)\le C\,\frac MN\,p^{(3s)}(z)$, by (P4) with $u=z$. Since $M/N\le\delta'^2\le\delta'm^{-1/2}$ and $p^{(3s)}\le2^{d/2}p^{(6s)}$, both are at most $C\delta'm^{-1/2}p^{(6s)}(z)$.

\emph{Conclusion.} On the intersection of these events, $|S_1|\le C\delta'm^{-1/2}p^{(6s)}(z)$ and
\[
\begin{aligned}
A&\le\mathbb E|Y_j|+C\,\delta'm^{-1/2}p^{(6s)}(z)\le C\,m^{-1/2}\big(p^{(2s)}(0)p^{(3s)}(z)\big)^{1/2}+C\,\delta'm^{-1/2}p^{(6s)}(z)\\
&\le C\,m^{-1/2}\,p^{(6s)}(z),
\end{aligned}
\]
again by (P5). With $\delta\le6\delta'$ and $p^{(6s)}=p_{6m}$ this is (ii). There are at most $N^2\delta'^{-2}\le N^3$ triples $(b,i,m)$ and two events for each, so (ii) fails with probability at most $4N^{-3}$; together with (i) and (iii) the total is at most $7N^{-3}$. The last two statements follow from $p_{6m}\le m_{6m}$ and from (iii) with $9\delta'\le1$.
\end{proof}

\begin{theorem}[The kernel of the sampled walk, uniformly in the number of steps]\label{thm:kernel}
In the situation of Lemma~\ref{lem:consistency} let $m_s=p^{(s\xi^2)}(0)$ for real $s>0$, let
\[
D_k=\max_{a,i}\big|N(P^k)_{ai}-p_k(x_a-x_i)\big|,
\]
and let $C_1=24+3\cdot2^{d+3}C$ with the constant $C$ of the lemma. On the event of the lemma, for all integers $k$ with $1\le k\le(C_1\delta')^{-2}$,
\[
D_k\le C_1\,\delta'\,\sqrt k\;m_k,\qquad\text{in particular}\qquad\big|\Tr P^k-m_k\big|\le C_1\,\delta'\,\sqrt k\;m_k .
\]
\end{theorem}

\begin{proof}
By (P2) and (P1), $s\mapsto m_s$ is non-increasing and $m_s\le2^{d/2}m_{2s}$.

\emph{Step 1: telescoping.} Let $R$ be the restriction of functions on $\M$ to the points, and write $(Pg)(b)=\sum_jP_{bj}\,g(x_j)$. For every $g$ and $k\ge1$,
\[
P^kRg-R\mathcal P^kg=\sum_{l=0}^{k-1}P^l\,(PR-R\mathcal P)\,\mathcal P^{k-1-l}g .
\]
For $g_i=p(\cdot-x_i)$ one has $\mathcal P^ng_i=p_{n+1}(\cdot-x_i)$ and $(PR-R\mathcal P)\,p_m(\cdot-x_i)=h^{(i)}_m$. With $m=k-l$,
\begin{equation}\label{eq:kernel-telescope}
(P^kRg_i)(a)=p_{k+1}(x_a-x_i)+\sum_{m=1}^k\big(P^{k-m}h^{(i)}_m\big)(a).
\end{equation}
\emph{Step 2: the last step of the walk.} As in Step~2 of the proof of Theorem~\ref{thm:jointpoly}, $NP_{bi}=p(x_b-x_i)(1+\theta_{bi})$ with $|\theta_{bi}|\le4\delta$ for $\delta\le\frac18$. Since $N(P^{k+1})_{ai}=\sum_b(P^k)_{ab}\,NP_{bi}$ and the entries of $P^k$ are nonnegative, $|N(P^{k+1})_{ai}-(P^kRg_i)(a)|\le4\delta\,(P^kRg_i)(a)$. With \eqref{eq:kernel-telescope} and $p_{k+1}\le m_{k+1}$,
\begin{equation}\label{eq:kernel-rec}
D_{k+1}\le4\delta\,m_{k+1}+(1+4\delta)\,S_k,\qquad S_k=\max_{a,i}\Big|\sum_{m=1}^k\big(P^{k-m}h^{(i)}_m\big)(a)\Big| .
\end{equation}
\emph{Step 3: two bounds for one term.} $P^l$ is row-stochastic, hence $|(P^lh)(a)|\le\max_b|h(b)|$; and for $l\ge1$, $|(P^lh)(a)|\le\max_bN(P^l)_{ab}\cdot\frac1N\sum_b|h(b)|\le(m_l+D_l)\,\frac1N\sum_b|h(b)|$. By the lemma, $\max_b|h^{(i)}_m(b)|\le C\delta'm^{-1/2}m_{6m}$ and $\frac1N\sum_b|h^{(i)}_m(b)|\le2C\delta'm^{-1/2}$.

\emph{Step 4: induction.} We show $D_k\le C_1\delta'\sqrt k\,m_k$ for $1\le k\le(C_1\delta')^{-2}$. For $k=1$, $|NP_{ai}-p(x_a-x_i)|\le4\delta M\le24\,\delta'm_1$. Let the claim hold up to $k$, with $k+1\le(C_1\delta')^{-2}$; then $D_l\le m_l$ for $l\le k$, and all $m\le k$ are in the range of the lemma. In $S_k$ use the first bound of Step~3 for $m>k/2$ and the second for $m\le k/2$. For $m>k/2$, $6m\ge k+1$ and $m_{6m}\le m_{k+1}$. For $m\le k/2$, $k-m\ge k/2\ge(k+1)/4$, so $m_{k-m}+D_{k-m}\le2m_{(k+1)/4}\le2^{d+1}m_{k+1}$. Hence
\[
S_k\le C\,\delta'\,m_{k+1}\Big[\sum_{k/2<m\le k}m^{-1/2}+2^{d+2}\sum_{m\le k/2}m^{-1/2}\Big]\le2^{d+3}C\,\delta'\,\sqrt k\;m_{k+1},
\]
and by \eqref{eq:kernel-rec}, with $\delta\le6\delta'\le\frac18$, $D_{k+1}\le\big(24+\tfrac32\cdot2^{d+3}C\big)\delta'\sqrt{k+1}\,m_{k+1}\le C_1\delta'\sqrt{k+1}\,m_{k+1}$. Finally $\Tr P^k-m_k=\frac1N\sum_i\big(N(P^k)_{ii}-p_k(0)\big)$.
\end{proof}

\begin{corollary}[The joint regime, almost surely]\label{cor:kernel-trace}
Let $\M=\R^d/\mathbb Z^d$, $d\ge1$, with the uniform density and the periodised Gaussian kernel, and let $\xi_N\to0$ and $N\xi_N^{d+2}/\log N\to\infty$. Then almost surely, for all large $N$, all $t\in[\xi_N^2,1]$ and all integers $1\le k\le20\,\xi_N^{-2}$, with $\sigma=\xi_N^2/2$, $\tau=2t/\xi_N^2$ and $\delta'=\sqrt{M\log N/N}$,
\[
\Big|\frac{\Tr P^k}{\Tr e^{k\sigma\Delta}}-1\Big|\le C_1\,\delta'\sqrt k,\qquad\Big|\frac{\Tr e^{-t\Lc}-Ne^{-\tau}}{T_\sigma(t)}-1\Big|\le C_1\,\delta'\sqrt{8\tau+2\log N+2}+\frac{4M}{N^2},
\]
and both right-hand sides are at most $C\big(\log N\,(1+\xi_N^2\log N)/N\xi_N^{d+2}\big)^{1/2}$, which tends to zero. If moreover $N\xi_N^{d+6}/\log N\to\infty$, they are $o(\xi_N^2)$, and almost surely
\[
\frac{\Tr e^{-t_N\Lc}-Ne^{-\tau}}{\Tr e^{t_N\Delta}}=1+\frac{d(d+2)}{8\tau}+O(\tau^{-2})+O(\xi_N^2)+o(\xi_N^2)
\]
for every sequence $t_N\in[\xi_N^2,1]$: the rate $\xi^2/t$ of open problem~2 holds with the sharp coefficient, almost surely, in this regime.
\end{corollary}

\begin{proof}
The probabilities $7N^{-3}$ are summable, so almost surely the event of Lemma~\ref{lem:consistency} holds for all large $N$. Let $J=\lceil8\tau\rceil+\lceil2\log N\rceil$; then $J\le16\,\xi_N^{-2}+2\log N+2$ and $\delta'^2J\le CM\log N\,(\xi_N^{-2}+\log N)/N\to0$; for the second summand note that the assumption gives $M\le C\xi_N^{-d}\le C(N/\log N)^{d/(d+2)}$ for large $N$, hence $M(\log N)^2/N\le C(\log N)^{2-d/(d+2)}N^{-2/(d+2)}\to0$. Consequently $J$ and $20\,\xi_N^{-2}$ are at most $(C_1\delta')^{-2}$ for large $N$, and the conditions $M\log N\le N$ and $6\delta'\le\frac18$ hold. The first statement is the theorem with $m_k=\Tr e^{k\sigma\Delta}$. For the second, $\Tr e^{-t\Lc}-Ne^{-\tau}=\sum_{j\ge1}w_j\Tr P^j$ by Theorem~\ref{thm:poisson} and $T_\sigma(t)=\sum_{j\ge1}w_jm_j$ by Theorem~\ref{thm:rate}(i), with $w_j=e^{-\tau}\tau^j/j!$. For $j\le J$ use the theorem and $\sqrt j\le\sqrt J$. For $j>J$, where $j\ge2$: $\Tr P^j\le\Tr P^2$ by Proposition~\ref{prop:moments}(ii), $\Tr P^2\le m_2+D_2\le2m_2\le2M$ by the theorem, $m_j\le M$, and $\sum_{j>J}w_j\le N^{-2}$ as in Step~4 of the proof of Theorem~\ref{thm:jointpoly}. Finally $T_\sigma(t)\ge\sum_{j\ge1}w_j=1-e^{-\tau}\ge\frac34$, because $m_j\ge1$ and $\tau\ge2$. If $N\xi_N^{d+6}/\log N\to\infty$, then $\delta'^2J=o(\xi_N^4)$: for $\log N\le\xi_N^{-2}$ it is at most $CM\log N\,\xi_N^{-2}/N$, and otherwise at most $CM(\log N)^2/N$ with $\xi_N^2>1/\log N$ and $M\le C\xi_N^{-d}\le C(\log N)^{d/2}$. The last statement follows from the uniform form of \eqref{eq:rateexact} in Step~5 of the proof of Theorem~\ref{thm:jointpoly}.
\end{proof}

\begin{corollary}[The corrected estimator on the sampled graph]\label{cor:kernel-corrected}
In the situation of Corollary~\ref{cor:kernel-trace} let $T^N(t)=\Tr e^{-t\Lc}-Ne^{-\tau}$. Almost surely, for all large $N$ and all $t\in[\xi_N^2,1]$,
\[
\Big|\Big(T^N-\frac{\sigma t}2\,(T^N)''\Big)(t)-\Big(T_\sigma-\frac{\sigma t}2\,T_\sigma''\Big)(t)\Big|\le(1+3\tau)\,C_1\,\delta'\sqrt{8\tau+2\log N+4}\;T_\sigma(t)+(3+6\tau)\frac M{N^2},
\]
where the derivatives are taken with respect to $t$. Consequently the relative deviation of the corrected estimator of the sampled graph from $\Tr e^{t\Delta}$ is $O(\tau^{-2})+O\big(\tau\,\delta'\sqrt{\tau+\log N}\big)$.
\end{corollary}

\begin{proof}
Let $a_j=\Tr P^j$ and $e_j=a_j-m_j$ for $j\ge1$, and $w_j=e^{-\tau}\tau^j/j!$ for $j\ge0$, $w_{-1}=0$. Then $T^N-T_\sigma=\sum_{j\ge1}w_je_j$, a series that may be differentiated term by term because $|e_j|\le N+M$, and $\frac{\dd^2}{\dd\tau^2}w_j=w_{j-2}-2w_{j-1}+w_j$. Hence
\[
\frac{\dd^2}{\dd\tau^2}\big(T^N-T_\sigma\big)=\sum_{i\ge0}w_i\,e_{i+2}-2\sum_{i\ge0}w_i\,e_{i+1}+\sum_{i\ge1}w_i\,e_i .
\]
Let $J$ be as in the proof of Corollary~\ref{cor:kernel-trace}; for large $N$, $J+2\le(C_1\delta')^{-2}$. For the indices up to $J+2$, $|e_j|\le C_1\delta'\sqrt{J+2}\,m_j$ by Theorem~\ref{thm:kernel}, and $\sum_{i\ge0}w_im_{i+2}\le\sum_{i\ge0}w_im_{i+1}\le e^{-\tau}M+T_\sigma\le\frac32T_\sigma$, because $m_j$ is non-increasing, $T_\sigma\ge w_1m_1=\tau e^{-\tau}M$ and $\tau\ge2$. The terms with $i\le J$ therefore contribute at most $\frac{11}2C_1\delta'\sqrt{J+2}\,T_\sigma$. For $i>J$ all indices are at least two, $|e_j|\le a_j+m_j\le3M$ as in the proof of Corollary~\ref{cor:kernel-trace}, and $\sum_{i>J}w_i\le N^{-2}$, which gives at most $12MN^{-2}$. Since $t=\sigma\tau$, $\frac{\sigma t}2\frac{\dd^2}{\dd t^2}=\frac\tau2\frac{\dd^2}{\dd\tau^2}$; together with $|T^N-T_\sigma|\le C_1\delta'\sqrt J\,T_\sigma+3MN^{-2}$ from the same proof this is the first statement, with $J+2\le8\tau+2\log N+4$. The second follows from Theorem~\ref{thm:corrected} and $T_\sigma\le C\,\Theta$.
\end{proof}

\begin{corollary}[The graph without self-loops]\label{cor:kernel-noloops}
Let $P$ be built as in Lemma~\ref{lem:consistency}, but with the diagonal of the kernel removed, that is, with $\tilde P_{bb}=0$ (the case $\rho=0$ of Proposition~\ref{prop:bias}). Then Lemma~\ref{lem:consistency} holds without change, $D_1=m_1$, and on the event of the lemma, with $C_1'=C_1+3\cdot2^{d/2}$,
\[
D_k\le C_1'\,\delta'\,\sqrt k\;m_k\qquad\text{for all integers }2\le k\le(C_1'\delta')^{-2}.
\]
In the situation of Corollary~\ref{cor:kernel-trace}, almost surely for all large $N$ and all $t\in[\xi_N^2,1]$,
\[
\Big|\Tr e^{-t\Lc}-Ne^{-\tau}-\big(T_\sigma(t)-\tau e^{-\tau}M\big)\Big|\le\Big(C_1'\,\delta'\sqrt{8\tau+2\log N+2}+\frac{4M}{N^2}\Big)\,T_\sigma(t),
\]
and $\tau e^{-\tau}M\le C\,\tau^{1+d/2}e^{-\tau}\,\Tr e^{t\Delta}$: the statements of Corollary~\ref{cor:kernel-trace} on the rate hold for the graph without self-loops as well.
\end{corollary}

\begin{proof}
\emph{The lemma.} In its proof the term $j=b$ of $S_1$ and of $A$ is absent, and in (i) the term $l=i$ is replaced by $-1/N$, since now $\epsilon_i=\frac1N\sum_{l\ne i}(p(x_i-x_l)-1)-\frac1N$; nothing else changes.

\emph{The recursion.} Now $NP_{bi}=p(x_b-x_i)(1+\theta_{bi})$ for $b\ne i$ and $NP_{ii}=0$, hence
\[
N(P^{k+1})_{ai}=\sum_b(P^k)_{ab}\,g_i(x_b)(1+\theta_{bi})-M\,(1+\theta_{ii})\,(P^k)_{ai},
\]
with $\theta_{ii}$ defined by the same formula as for $b\ne i$, $|\theta_{ii}|\le4\delta$. Since $N(P^k)_{ai}\le m_k+D_k$, \eqref{eq:kernel-rec} becomes
\[
D_{k+1}\le4\delta\,m_{k+1}+(1+4\delta)\,S_k+(1+4\delta)\,\frac MN\,(m_k+D_k).
\]
For $k=1$ the diagonal entry gives $D_1=m_1$, the other entries being at most $4\delta M$ apart. In the induction of Step~4 the hypothesis is used only in the form $D_l\le m_l$, which holds for $l=1$; the additional term is at most $3\frac MNm_k\le3\cdot2^{d/2}\delta'^2m_{k+1}$, because $M/N\le\delta'^2$ and $m_k\le2^{d/2}m_{2k}\le2^{d/2}m_{k+1}$. For $k=1$ one has $S_1\le C\delta'm_6\le C\delta'm_2$, so that $D_2\le C_1'\delta'\sqrt2\,m_2$, and the induction proceeds as before with $C_1'$ in place of $C_1$.

\emph{The trace.} Here $\Tr P=0$, so that $\Tr e^{-t\Lc}-Ne^{-\tau}=\sum_{j\ge2}w_j\Tr P^j$, to be compared with $\sum_{j\ge2}w_jm_j=T_\sigma(t)-\tau e^{-\tau}M$. The terms with $2\le j\le J$ and with $j>J$ are estimated as in the proof of Corollary~\ref{cor:kernel-trace}. Finally $M\le C\sigma^{-d/2}$ and $\Tr e^{t\Delta}\ge(4\pi t)^{-d/2}$ give $M\le C\tau^{d/2}\Tr e^{t\Delta}$.
\end{proof}

\begin{lemma}[The spectrum of the kernel matrix of the heat kernel]\label{lem:kernelmatrix}
Let $\M=\R^d/\mathbb Z^d$ with independent uniform points $x_1,\dots,x_N$, let $0=\lambda_1\le\lambda_2\le\dots$ be the eigenvalues of $-\Delta$ on $\M$ with multiplicities, and for an integer $k\ge1$ let $Q_k$ be the symmetric matrix with the entries $p_k(x_a-x_i)/N$ and the eigenvalues $\mu_1(Q_k)\ge\dots\ge\mu_N(Q_k)$. There are $C$ and $N_0$ depending only on $d$ such that for $N\ge N_0$ and $M\log N\le N$, with probability at least $1-4\cdot3^dN^{-5}$, for all $k\ge1$ and all $1\le j\le N$,
\[
\big|\mu_j(Q_k)-e^{-k\sigma\lambda_j}\big|\le C\,\sqrt{\frac{\log N}{N}}\;m_k .
\]
\end{lemma}

\begin{proof}
For $n\in\mathbb Z^d$ let $a_n=e^{-4\pi^2\sigma|n|^2}$, so that $p_k(x)=\sum_na_n^ke^{2\pi in\cdot x}$ and $m_k=\sum_na_n^k$. With $(v_n)_a=e^{2\pi in\cdot x_a}$, $Q_k=\frac1N\sum_na_n^kv_nv_n^*=XX^*$ for the Hilbert--Schmidt operator $X\colon\ell^2(\mathbb Z^d)\to\mathbb C^N$, $Xe_n=N^{-1/2}a_n^{k/2}v_n$. The non-zero eigenvalues of $XX^*$ and of $X^*X$ coincide with multiplicities, and $X^*X=A^{1/2}GA^{1/2}$ with $A=\operatorname{diag}(a_n^k)$ and $G_{nn'}=\frac1N\sum_ae^{2\pi i(n'-n)\cdot x_a}=:\hat g(n'-n)$, $G_{nn}=1$. The eigenvalues of $A$ in decreasing order are $e^{-k\sigma\lambda_j}$. By the min-max principle for compact self-adjoint operators,
\[
\big|\mu_j(Q_k)-e^{-k\sigma\lambda_j}\big|\le\big\|A^{1/2}(G-1)A^{1/2}\big\|\le\Big(\sum_{n\ne n'}a_n^ka_{n'}^k\,|\hat g(n'-n)|^2\Big)^{1/2}.
\]
For $m\ne0$ the real and the imaginary part of $\hat g(m)$ are means of $N$ independent centred variables with values in $[-1,1]$; by Hoeffding's inequality and the union bound over the at most $(3N)^d$ vectors with $0<|m|_\infty\le N$, with probability at least $1-4\cdot3^dN^{-5}$ all of them satisfy $|\hat g(m)|\le\gamma:=2\sqrt{2(d+5)\log N/N}$. The pairs with $|n'-n|_\infty\le N$ then contribute at most $\gamma^2m_k^2$ to the sum. In the remaining pairs one of the two vectors has $|\cdot|_\infty>N/2$; with $|\hat g|\le1$ they contribute at most $2m_k\sum_{|n|_\infty>N/2}a_n^k\le2m_k\,e^{-k\sigma\pi^2N^2/2}\,m_{k/2}\le2^{1+d/2}m_k^2\,e^{-\sigma\pi^2N^2/2}$. Finally $M\log N\le N$ and $M\ge(2\pi\xi^2)^{-d/2}$ give $\sigma N^2\ge\frac1{4\pi}N^{2-2/d}(\log N)^{2/d}$, so that $e^{-\sigma\pi^2N^2/2}\le N^{-2}$ for $N\ge N_0$.
\end{proof}

\begin{theorem}[The eigenvalues of the sampled graph on a flat torus, uniformly]\label{thm:eigtorus}
In the situation of Theorem~\ref{thm:kernel} (flat torus, uniform density, periodised Gaussian kernel, graph with self-loops) let $1=s_1\ge s_2\ge\dots\ge s_N\ge0$ be the eigenvalues of $S$, let $\lambda^N_j=(1-s_j)/\sigma$ be the eigenvalues of $\Lc$, and let $\lambda_j$ be as in Lemma~\ref{lem:kernelmatrix}. On the intersection of the events of Lemmas~\ref{lem:consistency} and~\ref{lem:kernelmatrix}:
\begin{enumerate}[label=\textup{(\roman*)}]
\item for all integers $1\le k\le(C_1\delta')^{-2}$ and all $1\le j\le N$,
\[
\big|s_j^k-e^{-k\sigma\lambda_j}\big|\le C_3\,\delta'\,\sqrt k\;m_k ;
\]
\item for all $j\le N$ with $\lambda_j\le\xi^{-2}$ and $C_4\,\frac{\delta'}{\xi}\,(1\vee\lambda_j)^{(d-1)/2}\le1$,
\[
\Big|\lambda^N_j-\frac{1-e^{-\sigma\lambda_j}}{\sigma}\Big|\le C_4\,\frac{\delta'}{\xi}\,(1\vee\lambda_j)^{(d+1)/2},\qquad\text{and}\qquad0\le\lambda_j-\frac{1-e^{-\sigma\lambda_j}}{\sigma}\le\frac{\xi^2\lambda_j^2}{4}.
\]
\end{enumerate}
The constants $C_3$, $C_4$ depend only on $d$. In particular, if $\xi_N\to0$ and $N\xi_N^{d+2}/\log N\to\infty$, the two events hold almost surely for all large $N$.
\end{theorem}

\begin{proof}
$C$ denotes constants that depend only on $d$.

\emph{The matrix $S^k$.} $S$ is positive semidefinite, because $K_{ij}=p(x_i-x_j)$ is the Gram matrix of a function with positive Fourier coefficients and $S$ is obtained from $K$ by a congruence with a diagonal matrix. $S=D^{1/2}PD^{-1/2}$, hence $(S^k)_{ai}=(D_a/D_i)^{1/2}(P^k)_{ai}$, and $D_a=\sum_jW_{aj}=N\,u_ar_a$ in the notation of the proof of Lemma~\ref{lem:consistency}. For $\delta\le\frac18$, $u_ar_a\in[1-4\delta,1+4\delta]$ as in Step~2 of the proof of Theorem~\ref{thm:jointpoly}, so that $|(D_a/D_i)^{1/2}-1|\le8\delta$ and
\[
\big|N(S^k)_{ai}-p_k(x_a-x_i)\big|\le8\delta\,(m_k+D_k)+D_k\le16\,\delta\,m_k+D_k .
\]
$S^k-Q_k$ is a symmetric matrix whose entries are at most $(16\delta m_k+D_k)/N$ in absolute value; its operator norm is at most its largest row sum, $16\delta m_k+D_k$. The eigenvalues of $S^k$ in decreasing order are $s_j^k$. By Weyl's inequality, Theorem~\ref{thm:kernel} and Lemma~\ref{lem:kernelmatrix},
\[
\big|s_j^k-e^{-k\sigma\lambda_j}\big|\le16\,\delta\,m_k+D_k+C\sqrt{\frac{\log N}{N}}\,m_k\le C_3\,\delta'\sqrt k\,m_k ,
\]
which is (i); here $\delta\le6\delta'$ and $\sqrt{\log N/N}\le\delta'$.

\emph{From powers to eigenvalues.} For $x,y\in[0,1]$, $|x^k-y^k|\ge k\min(x,y)^{k-1}|x-y|$. Let $a=e^{-\sigma\lambda_j}$ and $\eta=C_3\delta'\sqrt k\,m_k$, and suppose $a^k\ge2\eta$. Then $s_j^k\ge\frac12a^k$, $\min(s_j,a)^{k-1}\ge(\frac12a^k)^{(k-1)/k}\ge\frac12a^k$, and by (i)
\begin{equation}\label{eq:eig-step}
\big|s_j-a\big|\le\frac{2\eta}{k\,a^k},\qquad\Big|\lambda^N_j-\frac{1-a}{\sigma}\Big|\le\frac{2\,\eta\,e^{k\sigma\lambda_j}}{k\sigma}.
\end{equation}
\emph{The choice of $k$.} Let $\lambda_*=1\vee\lambda_j\le\xi^{-2}$, so that $\sigma\lambda_*\le\frac12$, and $k=\lfloor1/\sigma\lambda_*\rfloor\ge1$. Then $\frac12\le k\sigma\lambda_*\le1$, hence $e^{k\sigma\lambda_j}\le e$, $\sqrt k\le(\sigma\lambda_*)^{-1/2}$, and $m_k=p^{(2k\sigma)}(0)\le C\lambda_*^{d/2}$ by (P2), because $2k\sigma\ge1/\lambda_*$ and $\lambda_*\ge1$. Therefore $\eta\le C\,\delta'\,\sigma^{-1/2}\lambda_*^{(d-1)/2}$, and the right-hand side of \eqref{eq:eig-step} is at most $4e\,\lambda_*\eta\le C\,\frac{\delta'}\xi\,\lambda_*^{(d+1)/2}$. The two conditions used, $a^k\ge e^{-1}\ge2\eta$ and $k\le1/\sigma\le(C_1\delta')^{-2}$, hold if $C_4\frac{\delta'}\xi\lambda_*^{(d-1)/2}\le1$ with a suitable $C_4$. The last inequality of (ii) is $0\le y-(1-e^{-y})\le\frac12y^2$ for $y=\sigma\lambda_j\ge0$.

\emph{Almost sure validity.} The probabilities $7N^{-3}$ and $4\cdot3^dN^{-5}$ are summable, and the conditions $M\log N\le N$, $6\delta'\le\frac18$ hold for large $N$ if $N\xi_N^{d}/\log N\to\infty$.
\end{proof}

\begin{corollary}[Eigenspaces on a flat torus]\label{cor:eigspace}
In the situation of Theorem~\ref{thm:eigtorus} let $\Lambda>0$ be an eigenvalue of $-\Delta$ with the multiplicity $r$, let $\mathcal N_\Lambda=\{n\in\mathbb Z^d:4\pi^2|n|^2=\Lambda\}$ and $J=\{j:\lambda_j=\Lambda\}$, let $g_\Lambda$ be the distance from $\Lambda$ to the rest of the spectrum of $-\Delta$, and let $w_n\in\mathbb C^N$ be the sampled Fourier mode, $(w_n)_a=N^{-1/2}e^{2\pi in\cdot x_a}$. There is a constant $C_5$ depending only on $d$ such that, on the events of the theorem, if $\Lambda\le\xi^{-2}$ and
\[
\vartheta_\Lambda:=C_5\,\frac{\delta'}{\xi}\,\frac{\Lambda^{(d+1)/2}}{g_\Lambda}\le1,
\]
then the eigenvalues $s_j$, $j\in J$, of $S$ are separated from the others, and the orthogonal projection $\Pi^N_\Lambda$ onto the span of the corresponding eigenvectors satisfies
\[
\big\|(1-\Pi^N_\Lambda)\,w_n\big\|\le\vartheta_\Lambda\qquad\text{for every }n\in\mathcal N_\Lambda .
\]
If moreover $r\gamma\le\frac12$ with $\gamma=2\sqrt{2(d+5)\log N/N}$, then $\|(1-\Pi^N_\Lambda)u\|\le\sqrt{2r}\,\vartheta_\Lambda$ for every unit vector $u$ in the span of the $w_n$, $n\in\mathcal N_\Lambda$, and this span has the dimension $r$ of the range of $\Pi^N_\Lambda$.
\end{corollary}

\begin{proof}
Let $k=\lfloor1/\sigma\Lambda\rfloor$ as in the proof of Theorem~\ref{thm:eigtorus}, so that $\frac12\le k\sigma\Lambda\le1$, and $a=e^{-k\sigma\Lambda}$; note $\Lambda\ge4\pi^2$.

\emph{Separation.} Let $\Lambda^-<\Lambda<\Lambda^+$ be the neighbouring eigenvalues of $-\Delta$; $\Lambda^+\le4\Lambda$, because $4\Lambda$ is an eigenvalue. Then $e^{-k\sigma\Lambda^-}-a\ge a\,k\sigma(\Lambda-\Lambda^-)\ge e^{-1}g_\Lambda/2\Lambda$ and $a-e^{-k\sigma\Lambda^+}\ge e^{-k\sigma\Lambda^+}k\sigma(\Lambda^+-\Lambda)\ge e^{-4}g_\Lambda/2\Lambda$. Hence the numbers $e^{-k\sigma\lambda_j}$ with $j\notin J$ are at a distance at least $G:=g_\Lambda/2e^4\Lambda$ from $a$. By Theorem~\ref{thm:eigtorus}(i) every $s_j^k$ is within $\eta_0=C_3\delta'\sqrt k\,m_k\le C\frac{\delta'}\xi\Lambda^{(d-1)/2}$ of $e^{-k\sigma\lambda_j}$. If $\eta_0\le G/2$, the eigenvalues $s_j^k$ of $S^k$ with $j\in J$ lie in $[a-\eta_0,a+\eta_0]$ and all others at a distance at least $G/2$ from $a$; since $s\mapsto s^k$ is increasing on $[0,1]$, $\Pi^N_\Lambda$ is the spectral projection of $S^k$ for the interval $(a-G/2,a+G/2)$.

\emph{An approximate eigenvector.} In the notation of the proof of Lemma~\ref{lem:kernelmatrix}, $Q_k=\sum_{n'}a_{n'}^kw_{n'}w_{n'}^*$, $\|w_{n'}\|=1$ and $w_{n'}^*w_n=\hat g(n-n')$, so that for $n\in\mathcal N_\Lambda$
\[
\big\|(Q_k-a)\,w_n\big\|=\Big\|\sum_{n'\ne n}a_{n'}^k\,\hat g(n-n')\,w_{n'}\Big\|\le\gamma\,m_k+\sum_{|n'|_\infty>N/2}a_{n'}^k\le C\,\delta'\,m_k ;
\]
here $|n|_\infty\le\sqrt\Lambda/2\pi\le1/2\pi\xi\le N/2$, because $M\log N\le N$ implies $\xi\ge1/\pi N$, so that $|n-n'|_\infty>N$ implies $|n'|_\infty>N/2$, and the last sum is estimated as in the lemma. With $\|S^k-Q_k\|\le16\delta m_k+D_k$ from the proof of the theorem, $\|(S^k-a)w_n\|\le\eta_1:=C\,\delta'\sqrt k\,m_k$.

\emph{Conclusion.} For a symmetric matrix $B$, a unit vector $w$ and the spectral projection $\Pi$ of $B$ for the interval $(a-\rho,a+\rho)$, the spectral theorem gives $\rho^2\|(1-\Pi)w\|^2\le\|(B-a)w\|^2$. With $B=S^k$ and $\rho=G/2$: $\|(1-\Pi^N_\Lambda)w_n\|\le2\eta_1/G\le C\frac{\delta'}\xi\Lambda^{(d+1)/2}/g_\Lambda$, and $\eta_0\le G/2$ holds under the same condition; this fixes $C_5$. For $u=\sum_{n\in\mathcal N_\Lambda}c_nw_n$ of norm one, $1=\|u\|^2\ge(1-(r-1)\gamma)\sum|c_n|^2$, hence $\sum|c_n|^2\le2$, the $w_n$ are linearly independent, and $\|(1-\Pi^N_\Lambda)u\|\le\sum|c_n|\,\vartheta_\Lambda\le\sqrt{2r}\,\vartheta_\Lambda$.
\end{proof}

\begin{corollary}[Delocalisation of the eigenvectors]\label{cor:deloc}
In the situation of Theorem~\ref{thm:eigtorus}, on its events, let $u\in\R^N$ be a unit eigenvector of $S$ for the eigenvalue $s_j$, with $\lambda_j\le\xi^{-2}$ and $C_6\frac{\delta'}\xi(1\vee\lambda_j)^{(d-1)/2}\le1$. Then
\[
\max_a|u_a|\le C_6\,(1\vee\lambda_j)^{d/4}\,N^{-1/2},
\]
with a constant $C_6\ge C_4$ depending only on $d$.
\end{corollary}

\begin{proof}
Let $k=\lfloor1/\sigma\lambda_*\rfloor$ with $\lambda_*=1\vee\lambda_j$, as in the proof of the theorem; there $s_j^k\ge\frac12e^{-k\sigma\lambda_j}\ge1/2e$ was shown, and $2k\le2/\sigma\le(C_1\delta')^{-2}$ if $C_6$ is large enough. Since $u=s_j^{-k}S^ku$ and $S^k$ is symmetric, the Cauchy--Schwarz inequality gives
\[
|u_a|\le2e\,\Big(\sum_i(S^k)_{ai}^2\Big)^{1/2}=2e\,\big((S^{2k})_{aa}\big)^{1/2}\le2e\,\Big(\frac{(1+16\delta)\,m_{2k}+D_{2k}}N\Big)^{1/2}\le C\,\Big(\frac{m_{2k}}N\Big)^{1/2},
\]
by the bound for the entries of $S^{2k}$ in the proof of the theorem, and $m_{2k}\le m_k\le C\lambda_*^{d/2}$.
\end{proof}

\begin{remark}[On the uniform estimate of the eigenvalues]\label{rem:eigtorus}
(i) Theorem~\ref{thm:eigtorus} compares the eigenvalues of the sampled graph with those of the population walk, $(1-e^{-\sigma\lambda_j})/\sigma=\lambda_j-\frac14\xi^2\lambda_j^2+O(\xi^4\lambda_j^3)$, in the order of their size and with multiplicities, uniformly in $j$ as long as $\frac{\delta'}\xi\lambda_j^{(d-1)/2}$ is small. The bias of the eigenvalues is thus of second order in $\xi$ with an explicit coefficient, and the random error is $C\frac{\delta'}\xi\lambda_j^{(d+1)/2}$ with $\frac{\delta'}\xi\le C(\log N/N\xi^{d+2})^{1/2}$; for a fixed eigenvalue this is the rate of Theorem~\ref{thm:rate2}, here without the assumption that the eigenvalue is simple. (ii) In the form \eqref{eq:eigrate2}: for $\lambda_j\ge1$ and $\xi\sqrt{\lambda_j}\le1$, $|\lambda^N_j-\lambda_j|\le\big(C\frac{\delta'}\xi\lambda_j^{(d-1)/2}+\frac14\xi^2\lambda_j\big)\lambda_j$, so that \eqref{eq:eigrate2} holds on the flat torus with $\eta_N=C\delta'\xi_N^{-d}\le C(\log N/N\xi_N^{3d})^{1/2}$, provided $N\xi_N^{3d}/\log N\to\infty$; in dimension one this is the condition $N\xi_N^3/\log N\to\infty$. For the trace this adds nothing to Corollary~\ref{cor:kernel-trace}. Open problem~1 is answered by this for the flat torus with the Gaussian profile in the stated regime, and not beyond: the factor $\lambda_j^{(d-1)/2}$ in the random error comes from the bound of the kernel in the supremum norm and is probably not sharp, and the eigenvectors are treated only in the Euclidean norm on the points: by Corollary~\ref{cor:eigspace} the sampled Fourier modes of an eigenvalue $\Lambda$ lie within $C\frac{\delta'}\xi\Lambda^{(d+1)/2}/g_\Lambda$ of the corresponding eigenspace of the graph; a bound for this distance in the supremum norm, as needed for the embedding, is not shown; what is shown is that the eigenvectors themselves are delocalised, $\max_a|u_a|\le C\lambda_j^{d/4}N^{-1/2}$ (Corollary~\ref{cor:deloc}). (iii) \emph{Numerical illustration.} For one sample on the flat two-torus with $\xi=0.1$ and $N=3000$ the fifth, ninth and twenty-first eigenvalue of $\Lc$ in increasing order, counting from zero, are $62.45$, $105.56$ and $156.34$; the population values $(1-e^{-\sigma\lambda_j})/\sigma$ are $65.23$, $109.19$ and $158.77$, and $\lambda_j$ itself is $78.96$, $157.91$ and $315.83$. In this run $\delta'/\xi=2.1$, so the hypotheses of the theorem do not hold; the numbers show that the eigenvalues follow the population values and not $\lambda_j$. Within a group of equal eigenvalues the eigenvalues of the graph split. For $\xi=0.05$ and three samples each, the mean of the relative deviation from the population value over the group is $-0.99$, $-1.06$, $-1.31$ times $M/N$ for $\lambda_j=39.5$, $79.0$, $157.9$ at $N=3000$ and $-1.09$, $-1.17$, $-1.50$ times $M/N$ at $N=1500$, and the standard deviation within the group is $0.6$ to $1.4$ times $N^{-1/2}$ at $N=3000$ and $0.8$ to $1.6$ times $N^{-1/2}$ at $N=1500$. This suggests a bias of the order $M/N$, as for the trace, and a random error of the relative order $N^{-1/2}$ for the single eigenvalue, much smaller than the bound of the theorem; neither is proven. A further test separates the two: for $\xi=0.1$ and $\lambda_j=39.5$, over $120$, $80$ and $40$ samples at $N=400$, $800$ and $1600$, the standard deviation of the mean over the group, multiplied by $N$, is $3.5$, $3.7$ and $3.3$, whereas that of the lowest eigenvalue of the group, multiplied by $N^{1/2}$, is $0.53$, $0.48$ and $0.51$. The mean over a group thus fluctuates at the order $1/N$ and the single eigenvalue at the order $N^{-1/2}$; the alternative that the group mean fluctuates at the order $N^{-1/2}$ is not compatible with these numbers. A possible reason is that the splitting of a group is, to first order, given by the block of $G-1$ on the corresponding Fourier modes, whose trace vanishes because $G_{nn}=1$. (iv) Lemma~\ref{lem:kernelmatrix}, Theorem~\ref{thm:eigtorus} and Corollaries~\ref{cor:eigspace}, \ref{cor:deloc} were written by one session of an AI system and have not yet been read by a second one, nor by a person.
\end{remark}

\begin{remark}[Scope and status of the kernel estimate]\label{rem:kernel}
(i) The argument is the comparison of two recursions in the supremum norm: the sampled walk is stable, because $P$ is row-stochastic, and consistent with the population walk by Lemma~\ref{lem:consistency}. It uses neither the walk $\tilde P$ with the exact normalisation nor an expansion in $\epsilon$, and no eigenvalues. Corollary~\ref{cor:kernel-trace} contains the second statement of Theorem~\ref{thm:jointpoly} and its first statement in a larger regime, in every dimension and almost surely: the exponent $2d+6$ is replaced by $d+6$ for the sharp rate, and the trace converges, with the rate $(\log N/N\xi_N^{d+2})^{1/2}$ up to a logarithm, as soon as $N\xi_N^{d+2}/\log N\to\infty$. Theorem~\ref{thm:kernel} and Corollaries~\ref{cor:kernel-trace}, \ref{cor:kernel-corrected} are stated for the graph with self-loops; Corollary~\ref{cor:kernel-noloops} treats the graph without them. Theorem~\ref{thm:jointpoly} and its proof are kept because the count of closed walks is of independent use (Propositions~\ref{prop:bias}, \ref{prop:variance}).

(ii) The estimate is not sharp for the trace. It bounds every entry of $P^k$, and the supremum over the entries is of the order $\delta'$; the mean over the diagonal is much smaller, of the order $M/N$ according to Propositions~\ref{prop:bias} and~\ref{prop:variance} and Remarks~\ref{rem:bias}(v), \ref{rem:variance}(v). The sharp rate under the condition $N\xi_N^{d+2}\to\infty$ of those remarks would need this averaging and remains a conjecture.

(iii) \emph{Numerical illustration.} For one sample on the flat two-torus with $\xi=0.05$ and $N=4000$, where $\delta'=0.36$, the ratio $D_k/m_k$ is $0.51$, $0.39$, $0.32$, $0.24$, $0.21$, $0.20$, $0.20$ for $k=1,2,4,\dots,64$, and $\max_{b,i}|h^{(i)}_m(b)|\,m^{1/2}/p_{6m}(x_b-x_i)$ divided by $\delta'$ is $1.45$, $1.66$, $1.71$, $1.58$, $1.01$, $0.46$, $0.12$ for $m=1,2,4,\dots,64$; for $N=2000$ the numbers are similar. In these runs $\delta'$ is not small and the hypotheses of the lemma do not hold; they show the shape of the two bounds and not their validity.

(iv) \emph{Provenance.} Lemmas~\ref{lem:pgauss} and~\ref{lem:consistency} and Theorem~\ref{thm:kernel} with Corollary~\ref{cor:kernel-trace} were written on one day by two sessions of an AI system, each reading the part of the other line by line; of Corollary~\ref{cor:kernel-corrected} the second session has read the statement and a summary of the proof, not the text above, and Corollary~\ref{cor:kernel-noloops} has been read by both. No person has read them. Of all statements of this appendix these are the ones for which an independent check matters most.
\end{remark}

